\section{Constrained variational problems}

An integral constraint restricts an entire function rather than a finite list of coordinates. It can prescribe a total mass, average value, enclosed area, or total length while leaving the detailed shape free. For example, require
\begin{equation}
\int_a^b u(x)\,\mathrm{d}x=A
\end{equation}
while extremizing
\begin{equation}
J[u]=\int_a^b L(x,u,u')\,\mathrm{d}x.
\end{equation}
Introduce a constant multiplier $\lambda$ and define
\begin{equation}
\widetilde J[u]=\int_a^b\bigl(L(x,u,u')+\lambda u\bigr)\,\mathrm{d}x-\lambda A.
\end{equation}
Varying with respect to $\lambda$ restores the constraint, while varying with respect to $u$ determines the shape. The constant term does not affect the variation with respect to $u$, and the constrained problem becomes an ordinary Euler--Lagrange problem:
\begin{equation}
\frac{\partial L}{\partial u}+\lambda
-\frac{\mathrm{d}}{\mathrm{d}x}\left(\frac{\partial L}{\partial u'}\right)=0.
\end{equation}

\begin{example}[Shortest curve with prescribed area]
For the shortest curve between fixed endpoints subject to $\int u\,\mathrm{d}x=A$, take
\begin{equation}
L=\sqrt{1+u'^2}.
\end{equation}
The Euler--Lagrange equation is
\begin{equation}
\lambda-\frac{\mathrm{d}}{\mathrm{d}x}
\left(\frac{u'}{\sqrt{1+u'^2}}\right)=0.
\end{equation}
Integrating once gives
\begin{equation}
\frac{u'}{\sqrt{1+u'^2}}=\lambda x+c.
\end{equation}
After a second integration, the extremal satisfies an equation of the form
\begin{equation}
(\lambda x+c)^2+(\lambda u-d)^2=1,
\end{equation}
so the constrained shortest curve is an arc of a circle. The constants are fixed by the endpoint and area conditions. Equivalently, differentiating the first integral shows that the curvature is constant. The area requirement therefore changes the straight-line solution of the unconstrained problem into a constant-curvature curve.
\end{example}

The multiplier method for integral constraints is the function-space analogue of finite-dimensional Lagrange multipliers. The constraint is enforced by varying with respect to $\lambda$, while the Euler--Lagrange equation follows from varying with respect to $u$. This division of roles is also the origin of the constant chemical potential in the normalized density problem discussed earlier.